\begin{frame}{Физическая модель человеческого ECAP}
\centering
\begin{tikzpicture}[node distance=12mm]
\node[process] (stim) {Импульс стимуляции\\ и геометрия};
\node[process,right=of stim] (exc) {Возбуждение волокон\\ и распространение};
\node[process,right=of exc] (reg) {Объёмная проводимость\\ и регистрация};
\draw[flow] (stim)--(exc);\draw[flow] (exc)--(reg);
\end{tikzpicture}
\SlideFigureCaption{Формирование и регистрация ECAP}
\SlideGap
\raggedright
\begin{itemize}
\item Модель рекрутирования задаёт физически допустимые характеристики возбуждения волокон. Оператор регистрации преобразует их в измеряемый сигнал.
\item Возбуждение волокон будет рассчитываться однократно.
\item Артефакт, шум и условия регистрации будут учитываться отдельно.
\item Проверка потребует независимых человеческих записей и известных параметров измерительного канала.
\end{itemize}
До независимой проверки по записям человека результат будет обозначаться как \textbf{ECAP-подобный сигнал}.
\SourceNote{Anaya et al., Neuromodulation, 2020. DOI: \href{https://doi.org/10.1111/ner.12965}{10.1111/ner.12965}.}
\end{frame}
